% =====================================================================
%  Выводы по лабораторной работе.
% =====================================================================

\labpart{Выводы}

В ходе выполнения лабораторной работы реализована система
автоматического реферирования текстовых документов по варианту 16
(классический реферат тремя методами -- позиционно-частотный метод,
TextRank, эмбеддинги -- и реферат в виде иерархического списка
ключевых слов, с перефразированием через LLM) как новый сценарий
использования поверх архитектуры, разработанной в ЛР1--ЛР2: добавлен
один новый сервис (\texttt{summarization-service}), три новые таблицы
БД и новая страница интерфейса, при этом ни один существующий контракт
системы (поиск, краулинг, определение языка) не был нарушен. Все три
метода подчиняются единой схеме (расчёт веса предложения, отбор $N$
предложений с наибольшим весом, возврат к исходному порядку), различаясь
только правилом расчёта веса; реферат-список ключевых слов строится
единообразно поверх весов терминов, уже посчитанных для отбора
предложений, и не зависит от выбранного метода.

На тестовой коллекции из 20 документов англоязычной Wikipedia TextRank
оказался почти в 14 раз быстрее метода на эмбеддингах (289{,}45~мс
против 4013{,}54~мс на документ, табл.~\ref{tab:metrics-summarization})
за счёт полностью локального вычисления без обращений к внешнему
API. Позиционно-частотный метод дал самую высокую степень сжатия
(70{,}4\,\%), метод на эмбеддингах -- самую низкую (65{,}9\,\%) --
разница объясняется систематическим смещением каждого метода к
определённому типу предложений (первые/длинные предложения у
Posd/Posp против семантически центральных/коротких у эмбеддингов). Сама
метрика степени сжатия в процессе работы над лабораторной была
пересчитана с отношения числа предложений (гарантированно одинакового
у всех методов при одинаковом \texttt{sentence\_count} и потому
непригодного для сравнения методов) на отношение длины текста в
символах, которое действительно от метода зависит -- показательный
пример того, что метрика сравнения обязана быть чувствительна именно к
тому, что отличает сравниваемые варианты. Разработанный интерфейс
поддерживает как разовое построение реферата одним или сразу всеми
методами (результаты выводятся друг под другом, а не в колонках), так
и пакетное тестирование на всей коллекции, с выгрузкой результата в
CSV/JSON и печатью отдельной страницы с содержимым (включая
перефразированный LLM текст, если он был построен) -- без интерфейса
приложения, что закрывает соответствующее требование методических
указаний. Список документов для выбора по всему приложению реализован
единым переиспользуемым комбобоксом с поиском на сервере
(\texttt{DocumentCombobox}), а не статичной загрузкой всей коллекции.
